\section*{第 \ref{ch:g9:elements-universe-earth} 章——宇宙和地球中的元素}

\begin{solution}{exo:g9:elements-universe-earth:1}
氢（约 \qty{74}{\percent}）和氦（约
\qty{25}{\percent}）。
\end{solution}

\begin{solution}{exo:g9:elements-universe-earth:2}
两者中都含有氧。
\end{solution}

\begin{solution}{exo:g9:elements-universe-earth:3}
存在于恒星内部，以及恒星爆炸之中。
\end{solution}

\begin{solution}{exo:g9:elements-universe-earth:4}
约 \qty{3.5}{\percent}。
\end{solution}

\begin{solution}{exo:g9:elements-universe-earth:5}
$46.6 + 27.7 = \qty{74.3}{\percent}$，约为四分之三。
\end{solution}

\begin{solution}{exo:g9:elements-universe-earth:6}
它们是质量很轻的气体，从年轻的地球逸入太空，
（且氦不形成能将其保留在岩石中的化合物）。
\end{solution}

\begin{solution}{exo:g9:elements-universe-earth:7}
$0.23 \times 70 = \qty{16.1}{kg}$.
\end{solution}

\begin{solution}{exo:g9:elements-universe-earth:8}
氧\omterm{def:g7:atoms-and-molecules:atom}{原子}的质量是氢\omterm{def:g7:atoms-and-molecules:atom}{原子}的 16 倍：氧\omterm{def:g7:atoms-and-molecules:atom}{原子}数目较少，
但总质量仍然更大。
\end{solution}

\begin{solution}{exo:g9:elements-universe-earth:9}
约 $0.081 \times 1000 = \qty{81}{kg}$ 的铝。
\end{solution}

\begin{solution}{exo:g9:elements-universe-earth:10}
氧和钙（此外，除地壳中的八种主要元素外，
氢、碳和磷也以少量存在于地壳中，
含量均较少）。
\end{solution}

\begin{solution}{exo:g9:elements-universe-earth:11}
$40 \times 1.67 \times 10^{-27} = \qty{6.68e-26}{kg}$;
$1.06 / 6.68 \times 10^{-26} \approx 1.6 \times 10^{25}$ 个钙\omterm{def:g7:atoms-and-molecules:atom}{原子}。
\end{solution}

\begin{solution}{exo:g9:elements-universe-earth:12}
\omterm{def:g5:irreversible-changes:chemical-change}{化学变化}既不会产生\omterm{def:g7:atoms-and-molecules:atom}{原子}，也不会消灭\omterm{def:g7:atoms-and-molecules:atom}{原子}：它们只是
从一种化合物转移到另一种（岩石、气体、糖类、生物体）。
地球上的碳\omterm{def:g7:atoms-and-molecules:atom}{原子}被反复利用。
\end{solution}

\begin{solution}{pb:g9:elements-universe-earth:1}
\textbf{1.} $0.61 \times 70 = \qty{42.7}{kg}$.

\textbf{2.} 碳 $0.23 \times 70 = \qty{16.1}{kg}$；氢
$0.10 \times 70 = \qty{7.0}{kg}$.

\textbf{3.} $61 + 23 + 10 = \qty{94}{\percent}$.

\textbf{4.} 水，\ce{H2O}。

\textbf{5.} $16 \times 1.67 \times 10^{-27} = \qty{2.67e-26}{kg}$.

\textbf{6.} 碳：$12 \times 1.67 \times 10^{-27} = \qty{2.00e-26}{kg}$；
氢：\qty{1.67e-27}{kg}。

\textbf{7.} 16.

\textbf{8.} \omterm{def:g9:inside-the-atom:nucleus}{电子}的质量约比\omterm{def:g9:inside-the-atom:proton}{核子}小 1836 倍：
因此\omterm{def:g9:inside-the-atom:nucleus}{电子}的质量不到总质量的千分之一。

\textbf{9.} $7.0 / 1.67 \times 10^{-27} \approx 4.2 \times 10^{27}$
个氢\omterm{def:g7:atoms-and-molecules:atom}{原子}。

\textbf{10.} $16.1 / 2.00 \times 10^{-26} \approx 8.0 \times 10^{26}$
个碳\omterm{def:g7:atoms-and-molecules:atom}{原子}。

\textbf{11.} $42.7 / 2.67 \times 10^{-26} \approx 1.60 \times 10^{27}$
个氧\omterm{def:g7:atoms-and-molecules:atom}{原子}。

\textbf{12.} 我们的计数 $1.60 \times 10^{27}$ 与估算值
$1.61 \times 10^{27}$ 相符：一个 \qty{70}{kg} 的人体约含 $1.6 \times
10^{27}$ 个氧\omterm{def:g7:atoms-and-molecules:atom}{原子}。
\end{solution}
